\ifdefined\inMainDoc\else
\documentclass[12pt,a4paper]{article}
% ================================================================
%  PRÉAMBULE COMMUN - ANNALES CORRIGÉES
%  Université de Kara - FaST - Licence Fondamentale MPC
%  Design optimisé impression noir et blanc
% ================================================================

% -- ENCODAGE & LANGUE --
\usepackage[utf8]{inputenc}
\usepackage[T1]{fontenc}
\usepackage[french]{babel}
\usepackage{lmodern}
\usepackage{microtype}

% -- MISE EN PAGE --
\usepackage[a4paper, top=2.8cm, bottom=2.5cm, left=2.5cm, right=2.5cm, headheight=14pt]{geometry}
\usepackage{setspace}
\setstretch{1.2}
\usepackage{parskip}
\setlength{\parindent}{1.2em}
\setlength{\parskip}{4pt}

% -- MATHS --
\usepackage{amsmath, amssymb, mathtools, amsthm}
\usepackage{mathrsfs}
\usepackage{dsfont}
\usepackage{bm}
\usepackage{cancel}
\usepackage{textcomp}
% Définition de \oiint (intégrale de surface fermée) sans package supplémentaire
\newcommand{\oiint}{\mathop{\iint\!\!\!\!\!\!\!\bigcirc}\limits}

% -- PHYSIQUE & CHIMIE --
\usepackage{physics}
\usepackage{siunitx}
% Lorsque physics est chargé, siunitx omet \qty. On le restaure ici :
\AtBeginDocument{\RenewCommandCopy\qty\SI}
\sisetup{locale = FR, detect-all, output-decimal-marker={,}}
% Commande de substitution pour les unités posant problème avec pdflatex
% (ohm, degreeCelsius, micro, angstrom, percent utilisent des caractères non-ASCII)
\newcommand{\qtyohm}[1]{\ensuremath{\num{#1}\,\Omega}}
\newcommand{\qtydegc}[1]{\ensuremath{\num{#1}\,{}^\circ\text{C}}}
\newcommand{\qtyangstrom}[1]{\ensuremath{\num{#1}\,\text{\AA}}}
\newcommand{\qtypercent}[1]{\ensuremath{\num{#1}\,\%}}
\newcommand{\qtyum}[1]{\ensuremath{\num{#1}\,\mu\text{m}}}
\usepackage[version=4]{mhchem}

% -- GRAPHISMES --
\usepackage{graphicx}
\usepackage{float}
\usepackage{caption}
\usepackage{subcaption}
\usepackage{wrapfig}
\usepackage{tikz}
\usetikzlibrary{calc, arrows.meta, patterns, shapes.geometric}
\usepackage{pgfplots}
\pgfplotsset{compat=1.18}

% -- TABLEAUX --
\usepackage{booktabs}
\usepackage{array}
\usepackage{tabularx}
\usepackage{multirow}

% -- LISTES --
\usepackage{enumitem}
\setlist[enumerate]{topsep=3pt, itemsep=2pt, parsep=0pt}
\setlist[itemize]{topsep=3pt, itemsep=2pt, parsep=0pt, label={.}}

% -- BOÎTES (noir et blanc) --
\usepackage{mdframed}
\usepackage{tcolorbox}
\tcbuselibrary{skins, breakable, theorems}

% Boîte EXERCICE : encadré avec filet épais, fond blanc
\newmdenv[
  linewidth=1.2pt,
  linecolor=black,
  roundcorner=3pt,
  skipabove=10pt,
  skipbelow=6pt,
  innertopmargin=8pt,
  innerbottommargin=8pt,
  innerleftmargin=10pt,
  innerrightmargin=10pt,
  backgroundcolor=white,
  frametitle={\bfseries\small Exercice},
  frametitlealignment=\raggedright,
  frametitleaboveskip=4pt,
  frametitlebelowskip=4pt,
]{boiteexercice}

% Boîte CORRIGÉ : fond gris très clair + filet fin
\newmdenv[
  linewidth=0.6pt,
  linecolor=black,
  leftline=true,
  rightline=false,
  topline=false,
  bottomline=false,
  linecolor=black,
  innerleftmargin=12pt,
  innerrightmargin=8pt,
  innertopmargin=6pt,
  innerbottommargin=6pt,
  backgroundcolor=gray!8,
  skipabove=4pt,
  skipbelow=8pt,
]{boitecorrige}

% Boîte RÉSUMÉ DE COURS : double encadré
\newmdenv[
  linewidth=0.8pt,
  linecolor=black,
  roundcorner=2pt,
  backgroundcolor=gray!12,
  innertopmargin=10pt,
  innerbottommargin=10pt,
  innerleftmargin=12pt,
  innerrightmargin=12pt,
  skipabove=12pt,
  skipbelow=8pt,
]{boiteresume}

% Boîte ATTENTION/REMARQUE : barre gauche épaisse
\newmdenv[
  leftline=true,
  rightline=false,
  topline=false,
  bottomline=false,
  linewidth=3pt,
  linecolor=black,
  backgroundcolor=gray!5,
  innerleftmargin=12pt,
  innerrightmargin=8pt,
  innertopmargin=5pt,
  innerbottommargin=5pt,
  skipabove=6pt,
  skipbelow=6pt,
]{boiteremarque}

% -- DIFFICULTÉ DES EXERCICES --
\newcommand{\facile}{$\star$}
\newcommand{\moyen}{$\star\!\star$}
\newcommand{\difficile}{$\star\!\star\!\star$}
\newcommand{\niveau}[1]{\hfill\textbf{#1}\hspace{2pt}}

% -- EN-TÊTES ET PIEDS DE PAGE --
\usepackage{fancyhdr}
\pagestyle{fancy}
\fancyhf{}
\renewcommand{\headrulewidth}{0.5pt}
\renewcommand{\footrulewidth}{0.3pt}

% -- HYPERLIENS --
\usepackage[hidelinks]{hyperref}
\hypersetup{
    colorlinks=false,
    pdftitle={Annale Corrigée - Licence MPC - Université de Kara}
}

% -- TITRES --
\usepackage{titlesec}
\titleformat{\section}{\Large\bfseries}{\thesection.}{0.8em}{}[\vspace{2pt}\hrule\vspace{4pt}]
\titleformat{\subsection}{\large\bfseries}{\thesubsection.}{0.7em}{}
\titleformat{\subsubsection}{\normalsize\bfseries}{\thesubsubsection.}{0.6em}{}

% -- THÉORÈMES --
\theoremstyle{definition}
\newtheorem{defn}{Définition}[section]
\newtheorem{prop}{Propriété}[section]
\newtheorem{thm}{Théorème}[section]
\newtheorem{corol}{Corollaire}[section]
\newtheorem{exemple}{Exemple}[section]

% -- COMMANDES MATHÉMATIQUES UTILES --
\newcommand{\vect}[1]{\overrightarrow{#1}}
\newcommand{\R}{\mathbb{R}}
\newcommand{\N}{\mathbb{N}}
\newcommand{\Z}{\mathbb{Z}}
\newcommand{\C}{\mathbb{C}}
\renewcommand{\d}{\mathrm{d}}
\newcommand{\deriv}[2]{\dfrac{\d #1}{\d #2}}
\newcommand{\dpart}[2]{\dfrac{\partial #1}{\partial #2}}
\newcommand{\rot}{\overrightarrow{\mathrm{rot}}\,}
\renewcommand{\div}{\mathrm{div}\,}
\renewcommand{\grad}{\overrightarrow{\mathrm{grad}}\,}
\newcommand{\e}{\mathrm{e}}
\newcommand{\ii}{\mathrm{i}}
\renewcommand{\Tr}{\mathrm{Tr}}

% Numérotation des équations par section
\numberwithin{equation}{section}

% -- RÉSULTATS ENCADRÉS --
\newcommand{\res}[1]{\[\boxed{#1}\]}

% -- REMARQUE INLINE --
\newcommand{\rmq}[1]{\begin{boiteremarque}\textbf{Remarque.} #1\end{boiteremarque}}

% -- MISE EN VALEUR DU RÉSULTAT FINAL --
\newcommand{\reponse}[1]{%
  \begin{center}
    \fbox{\begin{minipage}{0.92\textwidth}\centering #1\end{minipage}}
  \end{center}%
}


\lhead{\small\textit{Programmation Fortran - Licence 3}}
\rhead{\small\textit{FaST, Université de Kara}}
\cfoot{\thepage}

% Configuration pour le code Fortran
\usepackage{listings}
\lstset{
  language=[90]Fortran,
  basicstyle=\ttfamily\small,
  keywordstyle=\bfseries,
  commentstyle=\itshape,
  stringstyle=\ttfamily,
  numbers=left,
  numberstyle=\tiny,
  numbersep=8pt,
  frame=single,
  breaklines=true,
  tabsize=2,
  showstringspaces=false
}

\begin{document}
\fi

\begin{center}
  {\LARGE\bfseries PROGRAMMATION EN FORTRAN}\\[0.3cm]
  {\normalsize Licence 3 - Physique}\\[0.1cm]
  \rule{0.7\textwidth}{0.8pt}
\end{center}

\section*{Résumé du cours}

\begin{boiteresume}

\textbf{1. Structure d'un programme Fortran.}
\begin{lstlisting}
program nom_programme
  implicit none
  ! declarations
  ! instructions
end program nom_programme
\end{lstlisting}
La directive \texttt{implicit none} est indispensable : elle oblige la déclaration explicite de toutes les variables.

\textbf{2. Types de base.} \texttt{integer} (entier), \texttt{real} (flottant simple précision), \texttt{double precision} ou \texttt{real(8)} (double précision), \texttt{complex} (complexe), \texttt{logical} (booléen), \texttt{character} (chaîne).

\textbf{3. Structures de contrôle.}
\begin{itemize}[nosep]
  \item Conditionnelle : \texttt{if (cond) then ... else if ... else ... end if}
  \item Boucle \texttt{do} : \texttt{do i=debut, fin, pas ... end do}
  \item Boucle \texttt{do while} : \texttt{do while (cond) ... end do}
  \item Sortie de boucle : \texttt{exit} (sortie) et \texttt{cycle} (passage à l'itération suivante)
\end{itemize}

\textbf{4. Tableaux.} \texttt{real, dimension(n) :: tab} ou \texttt{real :: tab(n)}. Indices de 1 à $n$ par défaut. Allocation dynamique : \texttt{allocatable}, \texttt{allocate}, \texttt{deallocate}.

\textbf{5. Fonctions et sous-programmes.}
\begin{itemize}[nosep]
  \item \texttt{function nom(args) result(res)} : retourne une valeur.
  \item \texttt{subroutine nom(args)} : ne retourne pas de valeur, modifie les arguments.
  \item Appel d'une sous-routine : \texttt{call nom(args)}.
\end{itemize}

\textbf{6. Entrées/Sorties.} \texttt{read(*,*) x} (entrée standard), \texttt{write(*,*) x} ou \texttt{print*, x} (affichage).

Pour un fichier : \texttt{open(10, file='data.txt')}, \texttt{read(10,*) x}, \texttt{close(10)}.

\end{boiteresume}

\newpage
\section{Exercices fondamentaux}

\subsection*{Exercice 1 - Résolution d'un système linéaire 2×2 \niveau{\facile}}

\begin{boiteexercice}
Écrire un programme Fortran qui résout le système :
\[
\begin{cases}u_1 x + v_1 y = w_1 \\ u_2 x + v_2 y = w_2\end{cases}
\]
par la méthode de Cramer. Le programme doit vérifier l'existence d'une solution unique et afficher le résultat.
\end{boiteexercice}

\begin{boitecorrige}
\begin{lstlisting}
program systeme
  implicit none
  real :: u1, u2, v1, v2, w1, w2
  real :: delta, delta_x, delta_y, x, y

  ! Valeurs des coefficients
  u1 = 2.0;  u2 = 4.0
  v1 = 5.0;  v2 = 11.0
  w1 = 7.0;  w2 = 6.0

  ! Calcul du determinant principal
  delta = u1*v2 - u2*v1

  if (abs(delta) < 1.0e-6) then
    print *, "Le systeme n'a pas de solution unique."
    stop
  end if

  ! Determinants de Cramer
  delta_x = w1*v2 - w2*v1
  delta_y = u1*w2 - u2*w1

  x = delta_x / delta
  y = delta_y / delta

  print *, "Solution : x =", x, ", y =", y
end program systeme
\end{lstlisting}

\textbf{Résultat pour les valeurs données :} $\delta = 22-20=2$, $\delta_x = 77-30=47$... On trouve $x=47/2=23.5$? Vérifions : $u1=2$, $v1=5$, $u2=4$, $v2=11$, donc $\delta=2\times11-4\times5=22-20=2$. $\delta_x=7\times11-6\times5=77-30=47$, donc $x=23.5$. $\delta_y=2\times6-4\times7=12-28=-16$, donc $y=-8$.

La règle de Cramer : pour $Ax=b$, $x_i = \det(A_i)/\det(A)$ où $A_i$ est la matrice $A$ dont la $i$-ème colonne est remplacée par $b$.
\end{boitecorrige}

\subsection*{Exercice 2 - Racines d'un trinôme du second degré \niveau{\facile}}

\begin{boiteexercice}
Écrire un programme Fortran qui calcule les racines (réelles) de $ax^2+bx+c=0$. Le programme doit :
\begin{enumerate}
  \item Vérifier que $a\neq0$.
  \item Calculer le discriminant $\Delta = b^2-4ac$.
  \item Afficher les racines si elles existent, ou un message approprié.
\end{enumerate}
\end{boiteexercice}

\begin{boitecorrige}
\begin{lstlisting}
program trinome
  implicit none
  real, parameter :: epsilon = 1.0e-6
  real :: a, b, c, delta, x1, x2

  ! Valeurs des coefficients
  a = 1.0;  b = -5.0;  c = 6.0  ! x^2 - 5x + 6 = 0

  if (abs(a) < epsilon) stop "a doit etre non nul."

  delta = b*b - 4.0*a*c

  if (delta < -epsilon) then
    print *, "Pas de racine reelle (discriminant negatif)."
  else if (abs(delta) <= epsilon) then
    x1 = -b / (2.0*a)
    print *, "Racine double : x =", x1
  else
    x1 = (-b - sqrt(delta)) / (2.0*a)
    x2 = (-b + sqrt(delta)) / (2.0*a)
    print *, "Deux racines : x1 =", x1, ", x2 =", x2
  end if
end program trinome
\end{lstlisting}

Pour $x^2-5x+6=0$ : $\Delta = 25-24=1>0$, racines $x_1=2$ et $x_2=3$.
\end{boitecorrige}

\subsection*{Exercice 3 - Suite de Fibonacci et nombre d'or \niveau{\moyen}}

\begin{boiteexercice}
La suite de Fibonacci est définie par $u_0=u_1=1$ et $u_{n+1}=u_n+u_{n-1}$. Le rapport $u_{n+1}/u_n$ converge vers le nombre d'or $\varphi = (1+\sqrt{5})/2$.

Écrire un programme qui calcule les termes de la suite jusqu'à ce que $|u_{n+1}/u_n - u_n/u_{n-1}| < 10^{-5}$, puis affiche la valeur approchée et la valeur exacte du nombre d'or.
\end{boiteexercice}

\begin{boitecorrige}
\begin{lstlisting}
program nombre_dor
  implicit none
  real, parameter :: epsilon = 1.0e-5
  real :: u_prec, u_cour, somme
  real :: v_prec, v_cour, nombre_or
  integer :: n

  nombre_or = (1.0 + sqrt(5.0)) / 2.0

  u_prec = 1.0;  u_cour = 1.0
  v_prec = 1.0;  n = 1

  do
    somme  = u_cour + u_prec
    v_cour = (u_cour + u_prec) / u_cour
    u_prec = u_cour
    u_cour = somme
    n = n + 1
    if (abs(v_cour - v_prec) < epsilon) exit
    v_prec = v_cour
  end do

  write(*,'(A,F12.8)') "Approximation (n = ", real(n), " ) :", v_cour
  write(*,'(A,F12.8)') "Valeur exacte              :", nombre_or
end program nombre_dor
\end{lstlisting}

La convergence est atteinte après une vingtaine d'itérations. Le nombre d'or $\varphi \approx 1{,}61803399$.
\end{boitecorrige}

\subsection*{Exercice 4 - Crible d'Ératosthène \niveau{\moyen}}

\begin{boiteexercice}
Le crible d'Ératosthène détermine tous les nombres premiers jusqu'à $n$ :
\begin{enumerate}
  \item Créer un tableau de tous les entiers de 2 à $n$.
  \item Pour chaque entier $i$ de 2 à $\sqrt{n}$ : si $i$ n'a pas été rayé, déclarer $i$ premier et rayer tous ses multiples.
  \item Les entiers non rayés sont premiers.
\end{enumerate}
Écrire ce programme en Fortran pour $n=1000$.
\end{boiteexercice}

\begin{boitecorrige}
\begin{lstlisting}
program eratosthene
  implicit none
  integer, parameter    :: n = 1000
  integer, dimension(n) :: tab
  integer               :: imax, i, j, nb_premiers

  ! Initialisation
  do i = 2, n
    tab(i) = i
  end do
  tab(1) = 0  ! 1 n'est pas premier

  ! Crible
  imax = int(sqrt(real(n)))
  do i = 2, imax
    if (tab(i) /= 0) then
      j = i + i
      do while (j <= n)
        tab(j) = 0
        j = j + i
      end do
    end if
  end do

  ! Affichage et comptage
  nb_premiers = 0
  do i = 2, n
    if (tab(i) /= 0) then
      nb_premiers = nb_premiers + 1
      write(*,'(I5)', advance='no') tab(i)
    end if
  end do
  print *
  write(*,'(A,I4,A,I4)') "Premiers 2 a ", n, " : ", nb_premiers
end program eratosthene
\end{lstlisting}

Il y a 168 nombres premiers entre 2 et 1000.

La complexité de l'algorithme est $O(n\log\log n)$.
\end{boitecorrige}

\subsection*{Exercice 5 - Tri à bulles \niveau{\moyen}}

\begin{boiteexercice}
Écrire un sous-programme Fortran \texttt{tri\_bulles} qui trie un tableau de réels dans l'ordre croissant. L'algorithme compare des éléments consécutifs et les échange si nécessaire, jusqu'à ce qu'aucun échange ne soit effectué en un passage. Tester avec le tableau $[0.76, 0.38, 0.42, 0.91, 0.25, 0.13, 0.52, 0.69, 0.76, 0.98]$.
\end{boiteexercice}

\begin{boitecorrige}
\begin{lstlisting}
subroutine tri_bulles(tab, n)
  implicit none
  integer, intent(in)    :: n
  real, intent(inout)    :: tab(n)
  real                   :: temp
  logical                :: tri_termine
  integer                :: i

  do
    tri_termine = .true.
    do i = 1, n-1
      if (tab(i) > tab(i+1)) then
        temp      = tab(i)
        tab(i)    = tab(i+1)
        tab(i+1)  = temp
        tri_termine = .false.
      end if
    end do
    if (tri_termine) exit
  end do
end subroutine tri_bulles

program test_tri
  implicit none
  integer, parameter :: n = 10
  real :: tab(n) = [0.76, 0.38, 0.42, 0.91, 0.25, &
                    0.13, 0.52, 0.69, 0.76, 0.98]
  call tri_bulles(tab, n)
  print *, "Tableau trie :", tab
end program test_tri
\end{lstlisting}

Résultat attendu : $[0.13, 0.25, 0.38, 0.42, 0.52, 0.69, 0.76, 0.76, 0.91, 0.98]$.

Complexité en cas défavorable : $O(n^2)$.
\end{boitecorrige}

\subsection*{Exercice 6 - Intégration numérique \niveau{\difficile}}

\begin{boiteexercice}
Écrire un programme Fortran qui calcule numériquement $I = \int_0^\pi \sin(x)\,\d x$ en utilisant la méthode des trapèzes, pour un nombre de subdivisions $N$ donné. Comparer avec la valeur exacte $I=2$.

La règle des trapèzes : $I \approx \frac{h}{2}\left[f(a) + 2\sum_{k=1}^{N-1}f(x_k) + f(b)\right]$ avec $h=(b-a)/N$ et $x_k=a+kh$.
\end{boiteexercice}

\begin{boitecorrige}
\begin{lstlisting}
program integration_trapeze
  implicit none
  real, parameter :: pi = 3.14159265358979
  integer         :: N, k
  real            :: a, b, h, x, integrale, valeur_exacte

  a = 0.0;  b = pi
  valeur_exacte = 2.0  ! valeur exacte de l'integrale de sin sur [0,pi]

  print *, "  N       I_approx      Erreur"
  print *, "-----------"

  do N = 10, 1000, 100
    h = (b - a) / real(N)
    ! Termes extremes
    integrale = 0.5 * (sin(a) + sin(b))
    ! Termes intermediaires
    do k = 1, N-1
      x = a + k * h
      integrale = integrale + sin(x)
    end do
    integrale = integrale * h
    write(*,'(I5, 2F14.8)') N, integrale, abs(integrale - valeur_exacte)
  end do
end program integration_trapeze
\end{lstlisting}

Pour $N=10$ : erreur $\approx 0.0082$. Pour $N=1000$ : erreur $\approx 8\times10^{-7}$. La méthode des trapèzes est d'ordre 2 : l'erreur est en $O(h^2) = O(1/N^2)$.
\end{boitecorrige}


\newpage
\section{Exercices avancés de programmation Fortran}

\subsection*{Exercice 7 - Résolution d'un système linéaire 2×2 \niveau{\facile}}

\begin{boiteexercice}
Écrire un programme Fortran qui résout le système de 2 équations à 2 inconnues par la règle de Cramer :
\[
\begin{cases}u_1 x + v_1 y = w_1 \\ u_2 x + v_2 y = w_2\end{cases}
\]
Le programme doit gérer le cas où le déterminant est nul.
\end{boiteexercice}

\begin{boitecorrige}
\begin{lstlisting}
program systeme_2x2
  implicit none
  real :: u1, v1, w1, u2, v2, w2
  real :: delta, x, y

  u1=2.0; v1=5.0; w1=7.0
  u2=4.0; v2=11.0; w2=6.0

  delta = u1*v2 - u2*v1
  if (abs(delta) < 1.0e-6) then
    print *, "Determinant nul: pas de solution unique."
    stop
  end if

  x = (w1*v2 - w2*v1) / delta
  y = (u1*w2 - u2*w1) / delta
  print *, "Solution: x =", x, ", y =", y
end program systeme_2x2
\end{lstlisting}

Pour $u_1=2$, $v_1=5$, $w_1=7$, $u_2=4$, $v_2=11$, $w_2=6$ : $\Delta = 22-20=2$, $x = (77-30)/2 = 23{,}5$, $y = (12-28)/2 = -8$.
\end{boitecorrige}

\subsection*{Exercice 8 - Racines d'un trinôme \niveau{\facile}}

\begin{boiteexercice}
Écrire un programme Fortran pour calculer les racines d'un trinôme du second degré $ax^2 + bx + c$. Le programme doit vérifier que $a \neq 0$ et gérer les trois cas selon le signe du discriminant.
\end{boiteexercice}

\begin{boitecorrige}
\begin{lstlisting}
program trinome
  implicit none
  real, parameter :: eps = 1.0e-6
  real :: a, b, c, delta, x1, x2

  a = 3.0; b = 7.0; c = -11.0

  if (abs(a) < eps) stop "a doit etre non nul."

  delta = b*b - 4.0*a*c

  if (delta < -eps) then
    print *, "Pas de racine reelle."
  else if (abs(delta) < eps) then
    x1 = -b / (2.0*a)
    print *, "Racine double: x =", x1
  else
    x1 = (-b - sqrt(delta)) / (2.0*a)
    x2 = (-b + sqrt(delta)) / (2.0*a)
    print *, "Deux racines: x1 =", x1, ", x2 =", x2
  end if
end program trinome
\end{lstlisting}

Pour $a=3$, $b=7$, $c=-11$ : $\Delta = 49+132 = 181 > 0$. $x_1 = (-7-\sqrt{181})/6 \approx -3{,}42$, $x_2 \approx 1{,}08$.
\end{boitecorrige}

\subsection*{Exercice 9 - Suite de Fibonacci et nombre d'or \niveau{\moyen}}

\begin{boiteexercice}
Écrire un programme Fortran qui calcule le nombre d'or à partir de la suite de Fibonacci définie par $u_0=1$, $u_1=1$, $u_{n+1}=u_n+u_{n-1}$. Le rapport $u_{n+1}/u_n$ converge vers le nombre d'or $\varphi = (1+\sqrt{5})/2$. Arrêter quand la variation relative est inférieure à $10^{-5}$.
\end{boiteexercice}

\begin{boitecorrige}
\begin{lstlisting}
program nombre_dor
  implicit none
  real, parameter :: eps = 1.0e-5
  real :: u_prec, u_cour, somme
  real :: rapport_prec, rapport_cour
  real :: phi_exact
  integer :: n

  phi_exact = (1.0 + sqrt(5.0)) / 2.0
  u_prec = 1.0; u_cour = 1.0
  rapport_prec = 1.0
  n = 1

  do
    somme = u_cour + u_prec
    u_prec = u_cour
    u_cour = somme
    rapport_cour = u_cour / u_prec
    n = n + 1
    if (abs(rapport_cour - rapport_prec) < eps) exit
    rapport_prec = rapport_cour
  end do

  print *, "Iterations :", n
  print *, "Approx     :", rapport_cour
  print *, "Exact      :", phi_exact
  print *, "Err rel    :", abs(rapport_cour - phi_exact)/phi_exact
end program nombre_dor
\end{lstlisting}

La convergence est atteinte en environ 25 itérations.

Le nombre d'or est $\varphi \approx 1{,}61803398\ldots$
\end{boitecorrige}

\subsection*{Exercice 10 - Crible d'Ératosthène \niveau{\moyen}}

\begin{boiteexercice}
Écrire un programme Fortran qui détermine tous les nombres premiers dans l'intervalle $[2, N]$ à l'aide du crible d'Ératosthène. On rappelle l'algorithme : former un tableau de $2$ à $N$, puis pour chaque entier non rayé, rayer tous ses multiples. La recherche peut s'arrêter à $\sqrt{N}$.
\end{boiteexercice}

\begin{boitecorrige}
\begin{lstlisting}
program eratosthene
  implicit none
  integer, parameter :: N = 200
  integer :: tab(N), imax, i, j, count

  ! Initialisation
  do i = 2, N
    tab(i) = i
  end do
  tab(1) = 0  ! 1 n'est pas premier

  ! Crible
  imax = int(sqrt(real(N)))
  do i = 2, imax
    if (tab(i) /= 0) then
      j = i + i
      do while (j <= N)
        tab(j) = 0
        j = j + i
      end do
    end if
  end do

  ! Affichage
  count = 0
  print *, "Nombres premiers entre 2 et", N, ":"
  do i = 2, N
    if (tab(i) /= 0) then
      write(*,'(I5)', advance='no') tab(i)
      count = count + 1
    end if
  end do
  print *
  print *, "Nombre total de premiers:", count
end program eratosthene
\end{lstlisting}

Il y a 46 nombres premiers entre 2 et 200. Le dernier est 199 (premier).
\end{boitecorrige}

\subsection*{Exercice 11 - Méthode de Simpson pour l'intégration \niveau{\difficile}}

\begin{boiteexercice}
Écrire un programme Fortran qui calcule numériquement $I = \int_0^1 \e^{-x^2}\,\d x$ en utilisant la méthode de Simpson :
\[
I \approx \frac{h}{3}\left[f(a) + 4f(x_1) + 2f(x_2) + 4f(x_3) + \cdots + f(b)\right]
\]
où $h = (b-a)/N$ et $N$ est pair. Comparer avec la valeur de référence $I \approx 0{,}74682$.
\end{boiteexercice}

\begin{boitecorrige}
\begin{lstlisting}
program simpson
  implicit none
  real, parameter :: pi = 3.14159265358979
  integer :: N, k
  real    :: a, b, h, x, integrale, ref

  a = 0.0; b = 1.0
  ref = 0.74682413  ! valeur de reference

  print *, "   N         I_Simpson     Erreur"
  print *, "-----------------------------------"

  N = 10
  do while (N <= 10000)
    if (mod(N, 2) /= 0) N = N + 1  ! N doit etre pair
    h = (b - a) / real(N)

    ! Termes extremes
    integrale = exp(-a**2) + exp(-b**2)

    ! Termes impairs (coefficient 4)
    do k = 1, N-1, 2
      x = a + k * h
      integrale = integrale + 4.0 * exp(-x**2)
    end do

    ! Termes pairs (coefficient 2)
    do k = 2, N-2, 2
      x = a + k * h
      integrale = integrale + 2.0 * exp(-x**2)
    end do

    integrale = integrale * h / 3.0
    write(*,'(I6, 2F14.8)') N, integrale, abs(integrale - ref)
    N = N * 10
  end do
end program simpson
\end{lstlisting}

La méthode de Simpson est d'ordre 4 : l'erreur est en $O(h^4)$. Pour $N=10$, erreur $\approx 10^{-6}$ ; pour $N=100$, erreur $\approx 10^{-10}$.
\end{boitecorrige}

\subsection*{Exercice 12 - Produit de deux matrices \niveau{\moyen}}

\begin{boiteexercice}
Écrire un programme Fortran qui effectue le produit de deux matrices $A$ ($n\times m$) et $B$ ($m\times p$). Définir les dimensions à l'aide de paramètres. Tester avec $n=3$, $m=3$, $p=3$, $A = I_3$ (identité) et $B$ quelconque.
\end{boiteexercice}

\begin{boitecorrige}
\begin{lstlisting}
program produit_matrices
  implicit none
  integer, parameter :: n = 3, m = 3, p = 3
  real :: A(n,m), B(m,p), C(n,p)
  integer :: i, j, k

  ! Initialisation de A = identite
  A = 0.0
  do i = 1, min(n, m)
    A(i,i) = 1.0
  end do

  ! Initialisation de B (quelconque)
  data ((B(i,j), j=1,p), i=1,m) / &
       1.0, 2.0, 3.0, &
       4.0, 5.0, 6.0, &
       7.0, 8.0, 9.0/

  ! Produit matriciel
  do i = 1, n
    do j = 1, p
      C(i,j) = 0.0
      do k = 1, m
        C(i,j) = C(i,j) + A(i,k) * B(k,j)
      end do
    end do
  end do

  ! Affichage
  print *, "Matrice C = A * B :"
  do i = 1, n
    write(*,'(3F8.2)') (C(i,j), j=1,p)
  end do
end program produit_matrices
\end{lstlisting}

Puisque $A = I_3$, le résultat est $C = B$. Vérification : $C_{ij} = \sum_k \delta_{ik} B_{kj} = B_{ij}$.
\end{boitecorrige}


\newpage
\section{Exercices avancés : Fichiers, matrices, méthodes numériques}

\subsection*{Exercice 13 - Écrire à l'écran \niveau{\moyen}}

\begin{boiteexercice}
Écrire un programme Fortran qui demande à l'utilisateur le nom d'un fichier texte, ouvre ce fichier, puis l'affiche à l'écran ligne par ligne en faisant précéder chaque ligne de son numéro. Tester avec le fichier source du programme lui-même.
\end{boiteexercice}

\begin{boitecorrige}
\begin{lstlisting}
program ligne_fichier
  implicit none
  integer :: i, err
  character(len=30) :: fich, ncar
  character(len=250) :: ligne
  character(len=60)  :: chaine

  print*, 'Donner le nom du fichier : '
  read*, fich

  open(unit=10, file=fich)
  err = 0
  i   = 1

  do while (err == 0)
    read(10, fmt='(1a250)', iostat=err) ligne
    if (err /= 0) exit
    ! Nombre de caracteres utiles (sans blancs debut/fin)
    write(ncar, '(I3)') len(trim(adjustl(ligne)))
    ! Construction du format de sortie
    chaine = '(1I3, 1a, 1a' // trim(adjustl(ncar)) // ')'
    write(*, fmt=chaine) i, ' ', trim(adjustl(ligne))
    i = i + 1
  end do

  close(10)
end program ligne_fichier
\end{lstlisting}

Le programme utilise \texttt{iostat=err} pour détecter la fin de fichier (code $\neq 0$) sans provoquer d'erreur. Le format de sortie est construit dynamiquement à partir de la longueur effective de chaque ligne.
\end{boitecorrige}

\subsection*{Exercice 14 - Produit tensoriel \niveau{\difficile}}

\begin{boiteexercice}
Écrire un programme Fortran utilisant une fonction interne \texttt{rptens} qui renvoie le produit tensoriel de deux vecteurs réels de même taille $n$ donnés en entrée. La taille des vecteurs n'est pas un argument de la fonction. Tester avec $\vec{u}=(1,2,3)$ et $\vec{v}=(2,1,3)$.
\end{boiteexercice}

\begin{boitecorrige}
Le produit tensoriel de deux vecteurs $\vec{u}$ et $\vec{v}$ est la matrice $W_{ij} = u_i\,v_j$.
\begin{lstlisting}
program prog_rptens
  implicit none
  integer, parameter :: n = 3
  integer :: i
  real, dimension(n)    :: u, v
  real, dimension(n, n) :: w
  character(len=30) :: cn, chaine

  u = (/1., 2., 3./);  v = (/2., 1., 3./)

  ! --- appel fonction
  w = rptens(u, v)

  ! --- affichage
  write(cn, *) n
  chaine = '(' // trim(adjustl(cn)) // 'F10.5)'
  do i = 1, n
    write(*, fmt=chaine) w(i, :)
  end do

contains

  function rptens(u, v) result(w)
    implicit none
    real, dimension(:), intent(in) :: u, v
    real, dimension(size(u), size(v)) :: w
    integer :: i, j
    do i = 1, size(u)
      do j = 1, size(v)
        w(i, j) = u(i) * v(j)
      end do
    end do
  end function rptens

end program prog_rptens
\end{lstlisting}

Pour $\vec{u}=(1,2,3)$ et $\vec{v}=(2,1,3)$, le résultat est :
\[
W = \vec{u}\otimes\vec{v} =
\begin{pmatrix}2 & 1 & 3\\4 & 2 & 6\\6 & 3 & 9\end{pmatrix}.
\]
\end{boitecorrige}

\subsection*{Exercice 15 - Produit de matrices (avancé) \niveau{\difficile}}

\begin{boiteexercice}
Écrire un programme Fortran qui effectue le produit de deux matrices entières $A$ ($m_a \times n_a$) et $B$ ($m_b \times n_b$) via :
\begin{enumerate}
  \item Une sous-routine interne \texttt{matrix\_mult} dont les arguments d'entrée sont $A$, $B$ et leurs dimensions, et dont l'argument de sortie est $C$.
  \item Une sous-routine interne \texttt{matrix\_mult\_bis} qui fait la même chose sans passer les dimensions en argument (tableaux allocatables).
\end{enumerate}
Tester avec $m_a=2, n_a=3, m_b=3, n_b=2$.
\end{boiteexercice}

\begin{boitecorrige}
\begin{lstlisting}
program prog_mult
  implicit none
  integer, parameter :: ma=2, na=3, mb=3, nb=2
  integer :: i
  integer, dimension(ma, na) :: A
  integer, dimension(mb, nb) :: B
  integer, dimension(ma, nb) :: C

  A = reshape((/1, 6, 4, 5, 2, 8/), (/2, 3/))
  B = reshape((/6, 4, 2, 8, 4, 9/), (/3, 2/))

  ! --- appel subroutine avec dimensions en arguments
  call matrix_mult(A, B, C, ma, na, mb, nb)
  print*, 'calcul avec matrix_mult :'
  do i = 1, ma
    print*, C(i, :)
  end do

  ! --- appel subroutine sans dimensions en arguments
  call matrix_mult_bis(A, B, C)
  print*, 'calcul avec matrix_mult_bis :'
  do i = 1, ma
    print*, C(i, :)
  end do

contains

  subroutine matrix_mult(A, B, C, ma, na, mb, nb)
    implicit none
    integer, intent(in)  :: ma, na, mb, nb
    integer, dimension(ma, na), intent(in)  :: A
    integer, dimension(mb, nb), intent(in)  :: B
    integer, dimension(ma, nb), intent(out) :: C
    integer :: i, j, k, coeff
    do i = 1, ma
      do j = 1, nb
        coeff = 0
        do k = 1, na
          coeff = coeff + A(i, k) * B(k, j)
        end do
        C(i, j) = coeff
      end do
    end do
  end subroutine matrix_mult

  subroutine matrix_mult_bis(A, B, C)
    implicit none
    integer, dimension(:, :), intent(in)  :: A
    integer, dimension(:, :), intent(in)  :: B
    integer, dimension(:, :), intent(out) :: C
    integer :: i, j, k, coeff
    do i = 1, size(A, 1)
      do j = 1, size(B, 2)
        coeff = 0
        do k = 1, size(A, 2)
          coeff = coeff + A(i, k) * B(k, j)
        end do
        C(i, j) = coeff
      end do
    end do
  end subroutine matrix_mult_bis

end program prog_mult
\end{lstlisting}

Les deux sous-routines donnent le même résultat. La version \texttt{bis} est plus générale car elle utilise \texttt{size()} pour récupérer les dimensions.
\end{boitecorrige}

\subsection*{Exercice 16 - Résolution par dichotomie et Newton \niveau{\difficile}}

\begin{boiteexercice}
Soit la fonction $f(x) = x^3 - 4x - 2$ définie sur $\mathbb{R}$.
\begin{enumerate}
  \item Montrer que $f$ admet trois racines réelles $x_1, x_2, x_3$ dans les intervalles $[-2,-1]$, $[-1,0]$ et $[2,3]$.
  \item Rechercher par dichotomie la solution de $f(x) = 0$ dans l'intervalle $[2, 3]$. Effectuer 3 itérations en indiquant la précision à chaque étape.
  \item Reprendre la deuxième question en utilisant la méthode de Newton avec $x_0 = 4$ : $x_{i+1} = x_i - f(x_i)/f'(x_i)$.
\end{enumerate}
\end{boiteexercice}

\begin{boitecorrige}
\textbf{1. Vérification des racines.} On calcule les valeurs de $f$ aux bornes :
\begin{itemize}
  \item Sur $[-2,-1]$ : $f(-2) = -8+8-2 = -2 < 0$ et $f(-1) = -1+4-2 = 1 > 0$. Signe opposé $\Rightarrow$ racine.
  \item Sur $[-1,0]$ : $f(-1) = 1 > 0$ et $f(0) = -2 < 0$. Signe opposé $\Rightarrow$ racine.
  \item Sur $[2,3]$ : $f(2) = 8-8-2 = -2 < 0$ et $f(3) = 27-12-2 = 13 > 0$. Signe opposé $\Rightarrow$ racine.
\end{itemize}

\textbf{2. Méthode de dichotomie} sur $[a_0, b_0] = [2, 3]$ :

\begin{center}
\begin{tabular}{|c|c|c|c|c|c|}
\hline
$i$ & $a_i$ & $b_i$ & $x_i$ & $f(x_i)$ & $e_i = (b_i-a_i)/2$ \\
\hline
0 & 2 & 3 & 2.5 & 3.625 & 0.5 \\
1 & 2 & 2.5 & 2.25 & 0.3906 & 0.25 \\
2 & 2 & 2.25 & 2.125 & $-0.9043$ & 0.125 \\
3 & 2.125 & 2.25 & 2.188 & $-0.2825$ & 0.0625 \\
\hline
\end{tabular}
\end{center}
Après 3 itérations, $x^* \approx 2.188$.

\textbf{3. Méthode de Newton} avec $f'(x) = 3x^2-4$ et $x_0 = 4$ :

\begin{center}
\begin{tabular}{|c|c|c|c|}
\hline
$i$ & $x_i$ & $x_{i+1}$ & $|x_{i+1}-x_i|$ \\
\hline
0 & 4 & 2.95455 & 1.04545 \\
1 & 2.95455 & 2.41493 & 0.53962 \\
2 & 2.41493 & 2.23533 & 0.1796 \\
3 & 2.23533 & 2.21459 & 0.02074 \\
\hline
\end{tabular}
\end{center}
La solution est $x^* \approx 2.21459$. Newton converge plus rapidement que la dichotomie.
\end{boitecorrige}

\subsection*{Exercice 17 - Mouvement d'un objet en chute libre \niveau{\moyen}}

\begin{boiteexercice}
Un objet est lâché depuis une hauteur $h_0$ à partir du repos ($v_0 = 0$) sous l'influence de la gravité. On néglige la résistance de l'air. L'équation du mouvement est :
\[
h(t) = h_0 - \frac{1}{2}g\,t^2
\]
où $h_0$ est la hauteur initiale (en m), $g = \qty{9.81}{m.s^{-2}}$ et $h(t)$ la hauteur à l'instant $t$.

On veut déterminer le temps $t$ nécessaire pour que l'objet atteigne le sol ($h(t) = 0$).
\begin{enumerate}
  \item Vérifier s'il existe une racine de $f(t) = h(t) = 0$ dans l'intervalle $[0, 10]$ secondes.
  \item Combien d'itérations de dichotomie sont nécessaires pour une précision de $10^{-3}$ ?
  \item Appliquer la méthode de dichotomie pour $h_0 = 100\;\text{m}$ et $g = 9{,}81\;\text{m.s}^{-2}$ avec une tolérance de $0{,}05$.
\end{enumerate}
\end{boiteexercice}

\begin{boitecorrige}
\textbf{1.} On pose $f(t) = h_0 - \tfrac{1}{2}g\,t^2$. Pour $h_0 = 100\;\text{m}$ :
\[
f(0) = 100 > 0, \qquad f(10) = 100 - \tfrac{1}{2}\times9{,}81\times100 = -390{,}5 < 0.
\]
$f$ change de signe sur $[0,10]$ $\Rightarrow$ il existe bien une racine.

\textbf{2.} Le nombre d'itérations pour une précision $\varepsilon = 10^{-3}$ :
\[
n = \frac{\log(b-a) - \log(\varepsilon)}{\log 2} = \frac{\log(10) + 3}{\log 2} \approx \frac{4}{0{,}301} \approx 13{,}3.
\]
Il faut donc environ \textbf{14 itérations}.

\textbf{3.} Application de la dichotomie ($h_0=100$, tolérance $0{,}05$) :

\begin{center}
\begin{tabular}{|c|c|c|c|c|c|}
\hline
$i$ & $a_i$ & $x_i$ & $b_i$ & $f(x_i)$ & $e_i$ \\
\hline
1 & 0.0000 & 5.0000 & 10.000 & $-22.625$ & 5.0000 \\
2 & 0.0000 & 2.5000 & 5.0000 & 69.3437 & 2.5000 \\
3 & 2.5000 & 3.7500 & 5.0000 & 31.0234 & 1.2500 \\
4 & 3.7500 & 4.3750 & 5.0000 & 6.1152 & 0.6250 \\
5 & 4.3750 & 4.6875 & 5.0000 & $-7.7759$ & 0.3125 \\
6 & 4.3750 & 4.5313 & 4.6875 & $-0.7106$ & 0.1562 \\
7 & 4.3750 & 4.4531 & 4.5312 & 2.7323 & 0.0781 \\
8 & 4.4531 & 4.4922 & 4.5312 & 1.0183 & 0.0391 \\
\hline
\end{tabular}
\end{center}
Le temps de chute est $t \approx 4{,}4922\;\text{s}$. La valeur exacte est $t = \sqrt{2h_0/g} = \sqrt{200/9{,}81} \approx 4{,}515\;\text{s}$.
\end{boitecorrige}

\subsection*{Exercice 18 - Résolution d'une équation non-linéaire \niveau{\difficile}}

\begin{boiteexercice}
On se propose de résoudre numériquement l'équation $f(x) = 0$ par la méthode de la bissection (dichotomie) et par la méthode de Newton-Raphson.
\begin{enumerate}
  \item Écrire un programme Fortran qui résout l'équation par la méthode de la dichotomie.
  \item Écrire un programme Fortran qui résout l'équation par la méthode de Newton.
  \item Tester les deux programmes pour l'équation $x - \cos(x) = 0$ sur l'intervalle $[0{,}5\;;\;1{,}5]$ avec une précision $\varepsilon = 10^{-4}$.
  \item Comparer les résultats des deux méthodes.
\end{enumerate}
\end{boiteexercice}

\begin{boitecorrige}

\textbf{1. Programme Fortran pour la dichotomie :}
\begin{lstlisting}
program dichotomie
  implicit none
  real :: a, b, eps, x, fa, fx
  integer :: niter

  print*, 'Donner les deux bornes a et b :'
  read*, a, b
  print*, 'Donner la precision eps :'
  read*, eps

  fa = fct(a)
  niter = 0

  do while (abs(b - a) > eps)
    x = (a + b) / 2.0
    fx = fct(x)
    if (fa * fx < 0.0) then
      b = x
    else
      a  = x
      fa = fx
    end if
    niter = niter + 1
  end do

  print*, 'La solution est x =', x
  print*, 'Nombre d''iterations :', niter

contains

  function fct(x)
    implicit none
    real :: fct, x
    fct = x - cos(x)
  end function fct

end program dichotomie
\end{lstlisting}

\textbf{2. Programme Fortran pour Newton-Raphson :}
\begin{lstlisting}
program Newton
  implicit none
  real :: x, x0, eps, fx, fpx
  integer :: niter

  print*, 'Donner la valeur initiale x0 :'
  read*, x0
  print*, 'Donner la precision eps :'
  read*, eps

  x = x0
  niter = 0

  do while (abs(fct(x)) > eps)
    fx  = fct(x)
    fpx = fctp(x)
    if (fpx == 0.0) then
      print*, 'Derivee nulle, arret.'
      stop
    end if
    x = x - fx / fpx
    niter = niter + 1
  end do

  print*, 'La solution est x =', x
  print*, 'Nombre d''iterations :', niter

contains

  function fct(x)
    implicit none
    real :: fct, x
    fct = x - cos(x)
  end function fct

  function fctp(x)
    implicit none
    real :: fctp, x
    fctp = 1.0 + sin(x)
  end function fctp

end program Newton
\end{lstlisting}

\textbf{3 \& 4. Comparaison.} Sur $[0{,}5\,;\,1{,}5]$ avec $\varepsilon = 10^{-4}$ :

La dichotomie requiert environ $n = \lceil\log_2((1{,}5-0{,}5)/10^{-4})\rceil = 14$ itérations.
Newton, avec $x_0 = 1$, converge en \textbf{3--4 itérations} vers la solution $x^* \approx 0{,}7391$, qui vérifie $\cos(0{,}7391) \approx 0{,}7391$.

Newton est nettement plus rapide (convergence quadratique) mais nécessite de connaître la dérivée $f'(x) = 1 + \sin(x)$.
\end{boitecorrige}

\ifdefined\inMainDoc\else\end{document}\fi
